\section{Έλεγχοι, Κάλυψη Κώδικα και Όρια της Επαλήθευσης}
\label{sec:appendix_testing_coverage}


Η χρήση του \ENG{Codecov} γίνεται για την παρακολούθηση της κάλυψης των υφιστάμενων \ENG{tests}, καθώς και για την επισήμανση νέων αλλαγών
που δεν συνοδεύονται από αντίστοιχα σενάρια ελέγχου. Η εικόνα της κάλυψης συμπληρώνεται από τα αρχεία \texttt{\ENG{codecov.yml}}. Στα βασικά αποθετήρια,
τα \ENG{Codecov status checks} έχουν κυρίως ενημερωτικό χαρακτήρα,
ενώ και τα σχόλια στα \ENG{PR} είναι ρυθμισμένα ώστε να μην εμποδίζουν τη ροή ενσωμάτωσης αλλαγών.

Στο \texttt{\ENG{scorm-engine}}, η σουίτα του \ENG{Maven} συνδυάζεται με παραγωγή αναφορών μέσω \ENG{JaCoCo}
και εξαγωγή του \texttt{\ENG{target/site/jacoco/jacoco.xml}} ως \ENG{artifact} και, υπό προϋποθέσεις, προς το \ENG{Codecov}.
Στο \texttt{\ENG{scorm-engine-php-sdk}}, το \ENG{PHPUnit} παράγει το \texttt{\ENG{clover.xml}}, ενώ στον \texttt{\ENG{player}}
η διαδικασία δοκιμών βασίζεται στα \texttt{\ENG{tsx --test}} και \ENG{c8}, με παραγωγή του \texttt{\ENG{lcov.info}}.
Παράλληλα, και στα τρία αποθετήρια εντοπίζονται τοπικά \ENG{test files} που αντιστοιχούν στα σχετικά \ENG{workflows},
στοιχείο που υποδηλώνει τεκμηριωμένη, αν και όχι εξαντλητική, κάλυψη επιμέρους τμημάτων κώδικα.

Ως προς τις εκδόσεις, το μεγαλύτερο μέρος χρησιμοποιεί \ENG{semantic-release}, με συμβατική ανάλυση \ENG{commits},
παραγωγή \ENG{release notes} και ενημέρωση του \ENG{CHANGELOG}. Η υποδομή ακολουθεί διαφορετικό μοτίβο, όπου το \ENG{workflow}
υπολογίζει το επόμενο \ENG{tag}, το δημοσιεύει και δημιουργεί \ENG{GitHub release}.

\begin{table}[H]
\centering
\caption{Τρέχουσα εικόνα δοκιμών, κάλυψης κώδικα και \ENG{release support} στα βασικά αποθετήρια}
\label{tab:appendix_test_coverage_matrix}
\begingroup
\footnotesize
\setlength{\tabcolsep}{3pt}
\renewcommand{\arraystretch}{1.20}
\begin{adjustbox}{center,max width=\textwidth}
\begin{tabular}{|p{2.6cm}|p{4.1cm}|p{2.7cm}|p{2.8cm}|p{3.8cm}|}
\hline
Αποθετήριο & Τεκμηριωμένα \ENG{test artifacts} & Μορφή \ENG{coverage export} & \ENG{Codecov} ένταξη & Σχόλιο ερμηνείας \\
\hline
\texttt{\ENG{scorm-engine}} & \ENG{JUnit} tests για \ENG{manifest parsing}, \ENG{runtime adapters}, \ENG{state store}, \ENG{architecture rules} & \texttt{\ENG{jacoco.xml}} & Ναι, με \ENG{informational status} & Η κάλυψη συνδέεται με πραγματικό \ENG{unit-level corpus} \\
\hline
\texttt{\ENG{scorm-engine-php-sdk}} & \ENG{PHPUnit} tests για \ENG{APIs}, \ENG{transport}, \ENG{auth strategies}, \ENG{mappers} & \texttt{\ENG{clover.xml}} & Ναι, με \ENG{informational status} & Η \ENG{CI} ροή συμβαδίζει με τα τοπικά \ENG{tests} του \ENG{repository} \\
\hline
\texttt{\ENG{player}} & \ENG{Node tests} για \ENG{observed client}, \ENG{session state}, \ENG{terminate handler}, \ENG{view model builder} & \texttt{\ENG{lcov.info}} & Ναι, με \ENG{informational status} & Το \ENG{test suite} είναι εστιασμένο σε μικρές μονάδες και όχι σε πλήρες \ENG{browser integration} \\
\hline
\texttt{\ENG{example-lms-client}} & \ENG{Laravel test harness} και \ENG{placeholder} \texttt{\ENG{ExampleTest}} αρχεία & \texttt{\ENG{clover.xml}} στο \ENG{workflow} & Ναι, με \ENG{informational status} & Υπάρχει \ENG{pipeline readiness}, αλλά περιορισμένο τρέχον βάθος δοκιμών \\
\hline
\texttt{\ENG{central-docker-infrastructure}} & \ENG{validate-env}, \ENG{compose config}, \ENG{smoke checks}, \ENG{container diagnostics} & Δεν υπάρχει \ENG{line coverage artifact} & Όχι & Η επαλήθευση είναι λειτουργική και συνθετική, όχι γραμμοκεντρική \\
\hline
\end{tabular}
\end{adjustbox}
\endgroup

\end{table}
